\section{Discusión}\label{sec:discussion}

\subsection{Rendimiento frente a interpretabilidad}

Los resultados ordenan las tres familias. El EBM obtiene la mejor discriminación (ROC-AUC $0{,}7755 \pm 0{,}0077$), seguido de la regresión logística ($0{,}7515 \pm 0{,}0075$) y del árbol ($0{,}7372 \pm 0{,}0044$). La ventaja del EBM se repite en exactitud equilibrada ($0{,}6997$) y F1 ($0{,}5068$), mientras que el árbol conserva la mayor exactitud ($0{,}7793$) y la mayor precisión ($0{,}4805$). Esa combinación no designa un ganador absoluto. Si el uso previsto es priorizar un número reducido de revisiones con pocos falsos positivos, el árbol resulta defendible; si el objetivo es no dejar casos sin detectar, el EBM ofrece una sensibilidad notablemente mayor ($0{,}6405$ frente a $0{,}4890$ del árbol y $0{,}5955$ de la logística). La diferencia no puede atribuirse solo a la expresividad, ya que también cuentan la forma funcional del modelo, la penalización y el presupuesto de búsqueda.

La interpretabilidad tampoco es una sola dimensión. El árbol de referencia de la partición 1 se lee de un vistazo, con ocho reglas, tres variables y una jerarquía aplicable manualmente; entre las cinco particiones se utilizan de tres a cinco variables. Esa brevedad tiene un precio, porque solo una parte de los quince predictores participa en cada ajuste, las probabilidades son frecuencias de entrenamiento y la decisión cambia de forma abrupta con el umbral. La logística es compacta (quince a diecinueve coeficientes no nulos), pero exige manejar log-odds, categorías de referencia y contrastes por desviación típica; su lectura es más lenta que la del árbol aunque su forma sea global y aditiva. El EBM es el más expresivo, con quince funciones, cada una con su escala y su forma, algunas no monótonas. Inspeccionarlas cuesta más, pero permite ver dónde y con qué signo actúa cada variable, algo que el árbol solo insinúa. En suma, el árbol gana en acceso inmediato, la logística en compacidad y el EBM en fidelidad descriptiva; la elección depende de quién audite el modelo.

Con umbral $0{,}5$ los tres modelos superan la exactitud mayoritaria del $78{,}79$\,\%, pero detectan menos del $30$\,\% de los positivos. Con el umbral interno la sensibilidad sube entre 34,23 y 35,14 puntos porcentuales y la exactitud cae por debajo de esa referencia. No hay umbral universal sin costes, porque fijarlo decide qué fallo se tolera, no es una propiedad del modelo.

\subsection{De ocho a seis reglas}

Con umbral $0{,}25$, las hojas 6 y 7 del árbol predicen ambas clase 0 y las hojas 10 y 11 ambas clase 1. Fusionarlas produce dos reglas con probabilidades $0{,}0800$ y $0{,}5353$, de modo que el árbol puede presentarse como una lista de seis reglas con las mismas decisiones. El ahorro es real pero modesto, y se paga con probabilidades más gruesas, ya que se pierde la diferencia entre $0{,}0720$ y $0{,}1208$, y entre $0{,}4646$ y $0{,}6932$. Los otros dos desdoblamientos no pueden eliminarse sin perder las hojas positivas 4 y 14. Con umbral $0{,}5$ la lista podría ser aún más corta, pero la sensibilidad descendería, de modo que existe un compromiso entre simplicidad y rendimiento. Si el destino del modelo es una lista de comprobación, seis reglas son preferibles; si es la priorización fina de casos, conviene conservar las ocho.

\subsection{Mejoras y límites}

Las mejoras posibles se agrupan en tres frentes. En el árbol, la poda por coste-complejidad comentada, que nos permite crecer un árbol grande y recortar ramas con una penalización por hoja elegida por validación, podría sustituir la rejilla por un criterio sistemático. En la logística, los splines o las interacciones entre variables darían margen a efectos hoy forzados a ser lineales, a costa de perder la lectura de un coeficiente por variable. En el EBM, imponer monotonía en variables como los retrasos podría alinear las formas con el dominio y facilitaría su lectura.


Los límites son claros. Las cinco particiones proceden de entrenamientos solapados y no son réplicas independientes, de modo que su dispersión mide simplemente la sensibilidad al reparto. La familia se eligió con los mismos resultados que se comparan, lo que puede introducir optimismo; la exactitud equilibrada interna de $0{,}7007$ usada para seleccionar el umbral final no acredita una mejora externa sobre el $0{,}6997$ de la comparación. El test sin etiquetas no permite una comprobación independiente del rendimiento del modelo final; además, el carácter semisintético del ejercicio limita su extrapolación a una población real.
